% ============================================================
% Visualization Style Guide — ODE IVP Project
% Compile: tectonic -X compile visualization_guide.tex
% ============================================================
\documentclass[11pt]{article}
\usepackage[margin=2.3cm]{geometry}
\usepackage{amsmath,amssymb}
\usepackage{booktabs}
\usepackage{fontspec}
\usepackage{microtype}
\usepackage{xcolor}
\usepackage{listings}
\usepackage[hidelinks]{hyperref}
\usepackage{enumitem}

\setmainfont{Times New Roman}
\newcommand{\e}[1]{{\emojifont #1}}

\setlength{\parindent}{0pt}
\setlength{\parskip}{0.5em}
\renewcommand{\arraystretch}{1.3}

\lstset{basicstyle=\ttfamily\small,breaklines=true,frame=single,
        numbers=none,showstringspaces=false}

\begin{document}

\begin{center}
  {\LARGE\bfseries Visualization Style Guide}\\[0.4em]
  {\large ODE IVP Project}
\end{center}
\vspace{0.2em}
\hrule
\vspace{0.6em}

Applied to every figure that appears in the report or slides. The goal is not
aesthetics for their own sake: \textbf{a figure is a piece of evidence}, and
evidence must be legible, honest, and self-explanatory.

% ============================================================
\section*{Rules}
% ============================================================

\subsection*{1. Log-log for error-vs-step plots}
When plotting error against step size $h$ (or $\Delta t$), use \textbf{logarithmic
axes on both}. This is non-negotiable: the whole point is that order of
convergence appears as a straight line whose slope you can read and compare to
theory.

\begin{lstlisting}[language=Python]
ax.loglog(h_values, errors, "o-")
# add the theoretical slope for comparison
h_lin = np.logspace(-3, -1, 20)
ax.loglog(h_lin, C * h_lin**1.0, "k--", label="order 1")
\end{lstlisting}
Always overlay the theoretical slope. An order-1 plot without the reference line
is a claim without a benchmark.

\textbf{Confidence band on the fitted line.} The slope is the claim; report it
with its uncertainty. \texttt{np.polyfit} with \texttt{cov=True} returns the
slope, its standard error, and the covariance the band is built from:

\begin{lstlisting}[language=Python]
coef, cov = np.polyfit(np.log(h_values), np.log(errors), 1, cov=True)
slope, slope_err = coef[0], np.sqrt(cov[0, 0])
lx = np.linspace(np.log(h_values.min()), np.log(h_values.max()), 50)
line = coef[0]*lx + coef[1]
se = np.sqrt(cov[1,1] + 2*lx*cov[0,1] + lx**2*cov[0,0])
ax.loglog(np.exp(lx), np.exp(line), "k--", label="fit")
ax.fill_between(np.exp(lx), np.exp(line - 1.96*se),
                np.exp(line + 1.96*se), color="gray",
                alpha=0.3, label="95% band")
\end{lstlisting}

A band only means something if the data is repeatable. A deterministic
convergence table gives a \emph{narrow} band (a fit diagnostic); a Monte-Carlo
claim needs repeats across seeds, and the spread across repeats is the honest
band.

\subsection*{2. Always label axes and units}
\begin{itemize}[leftmargin=2em, itemsep=1pt]
  \item Axis label: quantity (unit). E.g.\ \texttt{step size h [s]},
    \texttt{global error [---]}.
  \item Title: state what the figure shows AND what conclusion it supports.
  \item Every figure gets a caption in the report that says what is shown and
    why it matters. ``Figure 3. Convergence'' is not a caption.
\end{itemize}

\subsection*{3. Use colour responsibly}
\begin{itemize}[leftmargin=2em, itemsep=1pt]
  \item \textbf{Never use rainbow / \texttt{jet} colormaps.} They are illegible in
    print and misleading (perceptually non-uniform). Use \texttt{viridis},
    \texttt{plasma}, or categorical palettes (\texttt{tab10}).
  \item Colourblind-safe pairs for overlaid curves: blue/orange, or use
    linestyles + markers instead of relying on colour alone.
  \item In error plots, use \texttt{o-} markers with lines, not bare markers.
\end{itemize}

\subsection*{4. Annotate the physics/math}
\begin{itemize}[leftmargin=2em, itemsep=1pt]
  \item Mark the interesting region (e.g.\ the $h$ at which stability kicks in)
    with a vertical line or shaded patch, and label it.
  \item If there is a round-off floor in a log-log plot, point at it: ``error
    flattens at $h \approx 10^{-6}$ due to double-precision round-off''. That is
    a \emph{finding}, not noise.
\end{itemize}

\subsection*{5. Reproducibility}
\begin{itemize}[leftmargin=2em, itemsep=1pt]
  \item Save every figure from \texttt{run\_all.py} with a fixed seed and fixed
    file name: \texttt{fig05\_order\_euler.png}. No timestamps or
    machine-specific junk in filenames.
  \item $\mathrm{DPI} \geq 150$ for the report; $\geq 200$ for slides.
  \item Every number quoted in the text must be reproducible from a figure or a
    printed line in the code output.
\end{itemize}

\subsection*{6. Time-series sanity}
\begin{itemize}[leftmargin=2em, itemsep=1pt]
  \item For ODE solutions over time, always show a reference (analytic or
    high-accuracy) when one exists.
  \item When comparing integrators, use the SAME time axis and SAME line styles,
    and a legend.
\end{itemize}

\subsection*{7. Cost--accuracy evidence (Learning Outcome D)}
\begin{itemize}[leftmargin=2em, itemsep=1pt]
  \item Plot \textbf{error versus cost}, not error versus nothing: the horizontal
    axis is the \emph{measured} cost (wall time, right-hand-side evaluations,
    Newton iterations, or accepted/rejected steps), the vertical axis the
    achieved accuracy against your reference solution.
  \item \textbf{Name the cost measure before you quote it.} ``Our implicit
    method is faster'' is not a result; ``at a final-time error of $10^{-6}$,
    implicit Euler took 0.8\,s and RK4 took 3.1\,s on the same machine'' is.
  \item \textbf{Match the target first.} Bring both methods to the same error
    level before comparing costs --- a fast method at a loose tolerance proves
    nothing. State the target in the caption.
  \item \textbf{Say what you did not count} (Jacobian assembly and
    factorisation, setup, compilation, plotting) and hold it constant across the
    methods being compared.
  \item \textbf{The negative result is still a result.} If the method you
    expected to win lost, plot it and say so --- that is a complete answer.
\end{itemize}

% ============================================================
\section*{Checklist before you submit a figure}
% ============================================================
\begin{itemize}[leftmargin=2.4em, itemsep=1pt]
  \item[$\square$] Log-log used where appropriate (error vs $h$, error vs $\Delta t$)
  \item[$\square$] Axes labelled with units
  \item[$\square$] Title states the conclusion, not just the content
  \item[$\square$] Theoretical slope / reference overlaid where applicable
  \item[$\square$] No rainbow colormap; colourblind-safe if colour-coded
  \item[$\square$] All curves have a legend
  \item[$\square$] Caption says what is shown AND why it matters
  \item[$\square$] If costs are compared: cost measure named, accuracy target matched, both methods on one error-vs-cost plot
  \item[$\square$] Reproducible from \texttt{run\_all.py} (fixed seed, fixed filename)
  \item[$\square$] $\mathrm{DPI} \geq 150$
\end{itemize}

% ============================================================
\section*{Anti-examples (what to avoid)}
% ============================================================
\begin{table}[h]
\centering
\small
\begin{tabular}{@{}p{3.2cm}p{5.5cm}p{2.5cm}@{}}
\toprule
Bad & Why & Good \\
\midrule
\texttt{plt.plot(h, err)} linear axes
  & order-1 error looks like a boring flat curve; you can't read the slope
  & log-log \\
\texttt{plt.imshow(data, cmap="jet")}
  & rainbow exaggerates differences, illegible
  & \texttt{viridis} \\
Error bar with only 1 run
  & MC/statistical claims need repeats
  & repeat with seeds, show spread \\
Figure with no caption in the report
  & reader can't tell what evidence it provides
  & one-sentence takeaway per figure \\
``Figure 4'' referenced but not discussed
  & decorative figure
  & either discuss it or delete it \\
\bottomrule
\end{tabular}
\end{table}

\end{document}